% Exemplo: um sudoku com o pacote sudoku e uma rolagem de dados com epsdice e números aleatórios do
% PGF (a semente fixa repete a mesma rolagem a cada compilação; mude-a para outra rolagem).
% Compilar: latexmk -pdf sudoku-dados.tex         Testado com: sudoku 1.0, epsdice 2.1, TeX Live 2024 ou mais recente
%%BEGIN:sudokupreambulo
\documentclass[12pt, a4paper]{article}
\usepackage[T1]{fontenc}
\usepackage[brazilian]{babel}
\usepackage{sudoku}                   % ambiente sudoku
\renewcommand*\sudokuformat[1]{\Large\sffamily#1}  % algarismos para a grade de 6 cm
\usepackage{epsdice}                  % \epsdice{n}: face de um dado
\usepackage{pgffor, pgfmath}          % \foreach e números aleatórios
%%END:sudokupreambulo

\begin{document}
\pagestyle{empty}

%%BEGIN:sudokugrade
\setlength{\sudokusize}{6cm}
\begin{center}
\begin{sudoku}
|5|3| | |7| | | | |.
|6| | |1|9|5| | | |.
| |9|8| | | | |6| |.
|8| | | |6| | | |3|.
|4| | |8| |3| | |1|.
|7| | | |2| | | |6|.
| |6| | | | |2|8| |.
| | | |4|1|9| | |5|.
| | | | |8| | |7|9|.
\end{sudoku}
\end{center}
%%END:sudokugrade

%%BEGIN:sudokudados
\pgfmathsetseed{2026}
Rolagem de cinco dados:
\foreach \i in {1,...,5} {%
  \pgfmathrandominteger{\face}{1}{6}%
  \epsdice{\face}\ %
}
%%END:sudokudados

\end{document}
